\chapter{Background and Literature Review}
\label{chap:background}

\chapterintro{%
This chapter develops the technical background for key-resource-aware candidate
evaluation and positions \QFlow\ within QKD networking, adaptive control, and
reconfigurable computing. It also defines a comparison taxonomy that separates
same-function hardware evidence from network simulations and component-level FPGA
studies.}

\section{Trusted-Relay QKD Networks}

QKD establishes secret material across a quantum channel and an authenticated
classical channel. Networking adds key management, trusted relay operation,
application interfaces, routing, monitoring, and recovery. Surveys describe QKD
networks as layered systems in which physical links generate secret material and
classical control software manages its use
\cite{cao2022qinternet,mehic2020networking}. In a trusted-relay architecture,
intermediate sites are trusted with the relevant key-handling operations. The trust
assumption differs from entanglement-repeater architectures and must be part of the
security policy.

Once link post-processing completes, the KMS stores classical key data and metadata.
For link $(i,j)$, useful control quantities include available key units
$K_{ij}(t)$, pool capacity $K_{ij}^{\max}$, effective generation rate
$\lambda_{ij}(t)$, consumption rate $C_{ij}(t)$, operational status, utilization,
and health indicators derived from QBER or other monitored values. The physical
fibre model may use attenuation coefficient $\alpha_{\mathrm{fiber}}$, but this
quantity is distinct from controller scoring coefficients.

\section{Key Pools and Resource Evolution}

A discrete-time key-pool balance can be written as
\begin{equation}
K_{ij}[n+1]
=\operatorname{clip}_{[0,K_{ij}^{\max}]}
\left(K_{ij}[n]+G_{ij}[n]-C_{ij}[n]-X_{ij}[n]\right),
\label{eq:key-pool-balance}
\end{equation}
where $G_{ij}[n]$ is newly admitted secret material, $C_{ij}[n]$ is consumed
material, and $X_{ij}[n]$ is material removed through expiry, revocation, alarms,
or administrative action. Every term must use the same unit, such as bits or fixed
key blocks. A normalized occupancy is
\begin{equation}
o_{ij}[n]=\frac{K_{ij}[n]}{K_{ij}^{\max}},
\qquad 0\leq o_{ij}[n]\leq1.
\label{eq:occupancy}
\end{equation}

Generation and consumption rates convey information that occupancy alone cannot.
Two links with equal current occupancy can have different future risk if one is
replenishing quickly and the other serves a high request load. A useful candidate
descriptor therefore combines present reserve with bounded rate, load, and health
information.

Freshness, cryptoperiod, or expiry is a KMS/security-policy property. A policy may
invalidate material after a configured time, assign a soft preference to newer
authorized batches, or revoke material after an alarm. These operations concern the
management of classical keys; they are not quantum-state evolution in memory.

\section{Hard Eligibility and Soft Preference}

For a request requiring $D[n]$ key units, a directed link can be represented by the
hard eligibility indicator
\begin{equation}
\chi_{ij}[n;D]
=\mathbf{1}\{\mathrm{status}_{ij}=\mathrm{UP}\}
 \mathbf{1}\{K_{ij}[n]\geq D[n]\}
 \mathbf{1}\{q_{ij}[n]\leq q_{\mathrm{abort}}\}
 \mathbf{1}\{\mathcal{M}_{ij}[n]\ \text{is valid}\},
\label{eq:link-eligibility-bg}
\end{equation}
where $\mathcal{M}_{ij}$ collects authenticated administrative metadata. Candidate
$\mathcal{R}_i$ is feasible when every constituent link is eligible:
\begin{equation}
v_i[n]=\prod_{(u,v)\in\mathcal{R}_i}\chi_{uv}[n;D].
\label{eq:route-eligibility-bg}
\end{equation}

Hard checks protect invariants. Soft preference is applied only after feasibility is
known and expresses engineering priorities among admissible candidates. This
separation avoids using a weighted score as a substitute for security or resource
constraints.

\section{Candidate-Route Evaluation}

Candidate-generation algorithms may use shortest paths, constrained path search,
multipath methods, optimization, or policy filters. \QFlow\ receives the resulting
bounded set $\mathcal{R}_{sd}$ rather than traversing the complete graph. Each
candidate descriptor contains fixed-width fields such as normalized distance,
bottleneck occupancy, replenishment condition, load or imbalance, operational risk,
hop count, and identifier. Candidate evaluation then becomes a regular reduction:
compute a score for every feasible descriptor and retain the best candidate under a
deterministic tie-break rule.

The distinction between generation and evaluation is important for both science and
implementation. Network-level papers may assess blocking, topology robustness, or
global path quality, while an evaluator implementation may assess arithmetic
exactness, cycles, post-route timing, and resources. The measurements are related
but not interchangeable.

\section{Multi-Objective Scoring}

Key-aware selection combines objectives that have different units. Each metric is
therefore normalized to a bounded dimensionless penalty before weighting. For
candidate $\mathcal{R}_i$, a four-term vector can be written as
\begin{equation}
\boldsymbol{\phi}_i[n]
=
\begin{bmatrix}
1-\widehat{o}_i[n]\\
1-\widehat{r}_i[n]\\
\widehat{u}_i[n]\\
\widehat{\rho}_i[n]
\end{bmatrix},
\label{eq:bg-penalty-vector}
\end{equation}
representing scarcity, recharge deficit, utilization or imbalance, and operational
risk. Scoring profile $\Pi_k$ supplies coefficient vector
$\boldsymbol{\alpha}^{(k)}$, giving local cost
\begin{equation}
c_i^{(k)}[n]
=\operatorname{sat}_{W_c}
\left(\sum_{r=1}^{4}\alpha_r^{(k)}\phi_{i,r}[n]\right).
\label{eq:bg-local-cost}
\end{equation}

Tchebycheff decomposition provides a scalar comparison while retaining an explicit
multi-objective interpretation \cite{zhang2007moead}. If
$\widehat{\boldsymbol{f}}_i^{(k)}$ is the normalized objective vector and
$\boldsymbol{z}^{\star}$ is an ideal point, the score is
\begin{equation}
g_i^{(k)}[n]
=\max_j\left\{
\lambda_{\mathrm{TCH},j}^{(k)}
\left|\widehat{f}_{i,j}^{(k)}[n]-z_j^{\star}\right|
\right\}.
\label{eq:bg-tchebycheff}
\end{equation}
The vector $\boldsymbol{\lambda}_{\mathrm{TCH}}^{(k)}$ contains objective weights;
it is distinct from link key-generation rate $\lambda_{ij}(t)$. The use of Pareto and
decomposition concepts is grounded in established multi-objective optimization
methods \cite{deb2002nsga2,zhang2007moead}.

\section{Offline Reinforcement Learning for Profile Selection}

Reinforcement learning (RL) maps observed state to actions that maximize expected
long-term reward. In tabular Q-learning, an offline training update is
\begin{equation}
Q(s_t,a_t)\leftarrow Q(s_t,a_t)+\eta_{\mathrm{RL}}
\left[r_t+\gamma_{\mathrm{RL}}\max_a Q(s_{t+1},a)-Q(s_t,a_t)\right],
\label{eq:q-learning}
\end{equation}
where $\eta_{\mathrm{RL}}$ is the learning rate and $\gamma_{\mathrm{RL}}$ is the
discount factor \cite{watkins1992qlearning,sutton2018rl}. In this thesis, the action
does not directly construct a route. It selects one of four fixed scoring profiles.
This reduces the deployed decision to a finite mapping from a quantized state to a
two-bit action.

Training and inference are separated. Training explores traces, updates Q-values,
and freezes a selected policy. Deployment stores only the policy mapping needed for
deterministic inference. The FPGA therefore contains no exploration mechanism and no
runtime Q-value update path. A minimum-dwell mechanism can suppress rapid profile
changes by holding the current action until a specified number of accepted decisions
has elapsed.

\section{Fixed-Point Arithmetic and FPGA Implementation}

Fixed-point arithmetic assigns an explicit integer and fractional width to every
quantity. For a signed Q$m.n$ value, representable values and rounding rules depend
on total width, signedness, scaling, saturation, and conversion policy. Bit-exact
verification requires these details to be identical in the software reference and
RTL. Saturation is preferable to silent wraparound for bounded metrics because it
preserves the intended ordering at numeric limits.

FPGA logic can implement quantizers, ROMs, adders, comparators, and finite-state
control with a deterministic clocked interface. Look-up tables (LUTs), registers,
occupied slices, block memories, and timing reports characterize the physical
mapping. Post-synthesis timing is an estimate before placement and routing;
post-route timing includes the implemented netlist and routing delays and is the
appropriate source for the maximum-frequency values used in the final comparison.

\section{CPU, FPGA, and ASIC Characteristics}

A CPU offers mature software tools, flexible control flow, and efficient handling of
irregular workloads. It is appropriate for topology management, candidate
generation, KMS/SDN orchestration, offline training, logging, and security
administration. Its execution time can depend on operating-system scheduling,
caches, and memory behavior unless a real-time environment is used.

An FPGA supports reconfigurable spatial datapaths and deterministic state machines.
It is attractive for an evolving fixed-point kernel because profiles and interfaces
can be revised by loading a new bitstream. A dedicated ASIC may ultimately provide
greater density, performance, and energy efficiency, especially at high volume, but
requires larger non-recurring engineering effort, signoff, fabrication, packaging,
test, and redesign if the contract changes. Platform comparison must therefore be
specific to development stage, volume, and measurement boundary.

\section{Literature Review}

\subsection{QKD routing and network control}

Bi et al. use virtual-network state and key resources for dynamic QKD routing
\cite{bi2023dynamic}. Chen et al. study hybrid-trusted path selection
\cite{chen2023hybrid}; Kiktenko et al. examine multiple non-overlapping paths
\cite{kiktenko2024nonoverlap}; and Akhtar et al. present fast and secure routing
algorithms for QKD networks \cite{akhtar2023fast}. Zhang et al. jointly optimize
routing, channels, key rate, and time slots \cite{zhang2024rckta}. These studies
show why hop count alone is insufficient, but their evaluated functions and evidence
boundaries differ from a fixed-point candidate-evaluation kernel.

RL-based work includes deep-RL resource assignment
\cite{sharma2023drl}, tabular Q-learning for resource-aware trusted-relay routing
\cite{hao2025qlearning}, and a graph-attention agent for flexible meshed topologies
\cite{johann2026adaptive}. These studies provide broader network or generalization
evidence than the present physical replay. They do not provide a same-input,
same-function FPGA timing baseline.

\subsection{Operational QKD systems}

MadQCI demonstrates heterogeneous SDN-QKD integration in production facilities
\cite{martin2024madqci}. Chen et al. report a carrier-grade quantum communication
network extending beyond 10,000~km \cite{chen2025carrier}. Such systems establish
the practical setting for KMS and controller integration. Their scale and end-to-end
objectives do not make their performance numbers directly comparable with a bounded
FPGA kernel.

\subsection{FPGA acceleration in QKD}

Reconfigurable hardware has accelerated multidimensional reconciliation encoding,
QKD post-processing, polar-code information reconciliation, error correction, and
synchronization
\cite{lu2023multidimensional,venkatachalam2023postprocessing,liao2025polar,
wei2025fptems,chandra2025synchronization}. These contributions demonstrate that FPGA
implementation is practical in QKD systems, but they process different functions
and use throughput or latency boundaries that cannot be converted into route-selection
speedups.

Hardware genetic-algorithm research provides relevant implementation ideas for
pipelined evaluation, selection, and specialization
\cite{shackleford2001pipelined,tang2004hardware,fernando2010ga,
guo2014automated,peker2018ga}. It remains component-level context because those
systems do not evaluate trusted-relay key-resource descriptors.

\subsection{Comparison taxonomy}

\begin{table}[H]
\centering
\caption{Comparison classes used throughout the thesis.}
\label{tab:comparison-classes}
\begin{tabularx}{\textwidth}{@{}p{0.10\textwidth}p{0.31\textwidth}X@{}}
\toprule
\textbf{Class} & \textbf{Definition} & \textbf{Valid comparison}\\
\midrule
A & Same function, board, interface, numeric contract, replay, and measurement
boundary & Direct timing, latency, resource, exactness, and controller-behavior
comparison\\
B & Same function and inputs on CPU/software and FPGA & Direct speed or energy
comparison only when both are measured under a matched protocol\\
C & QKD routing and network-control algorithms & Blocking, utilization, quality,
robustness, or generalization under compatible network experiments\\
D & Different QKD components, FPGA/RL hardware, deployments, or VLSI kernels &
Architectural and feasibility context; no raw QFlow speed ratio\\
\bottomrule
\end{tabularx}
\end{table}

Class A is supplied internally by the five same-board configurations. No Class B
experiment was completed. The QKD routing studies are Class C, while FPGA
post-processing, operational deployments, hardware genetic algorithms, and isolated
VLSI kernels are Class D. This taxonomy prevents a latency ratio between systems that
perform different computations.

\subsection{Reproducible literature search and verification}

The literature record was verified on 17 July 2026 from a predefined 50-row review
matrix. Each row was searched by exact title at an official publisher database or
the official arXiv record, including IEEE Xplore, ACM Digital Library, SpringerLink,
Nature, MDPI, Wiley/IET, Optica, Elsevier ScienceDirect, INFORMS, and Google Books
where appropriate. The audit checked bibliographic identity, platform, evidence
layer, metric scope, and whether decision latency was actually reported. One audit
identifier was retained per matrix row; sources whose exact identity could not be
verified were excluded, and explicit replacement sources were recorded as
replacements rather than silently treated as the same publication.

The resulting disposition contained 21 Class C and 29 Class D records, with no
identified Class A or Class B publication in the documented search. This is a
bounded finding from the recorded matrix and search date. Search queries, publisher
locations, verification outcomes, and the final workbook are preserved under
\path{literature/p1/} on the evidence branch identified in
Appendix~\ref{app:reproducibility}.

\subsection{Research position}

\begin{table}[H]
\centering
\caption{Representative related-work positioning. Direct numerical comparability
requires an equivalent function and measurement boundary.}
\label{tab:related-position}
\small
\begin{tabularx}{\textwidth}{@{}p{0.17\textwidth}p{0.22\textwidth}p{0.15\textwidth}X@{}}
\toprule
\textbf{Work/category} & \textbf{Function} & \textbf{Evidence} &
\textbf{Relationship to \QFlow}\\
\midrule
Hao et al. \cite{hao2025qlearning} & Resource-aware adaptive routing & Network
simulation & Class C RL counterpart; broader routing scope, no same-function FPGA\\
Johann et al. \cite{johann2026adaptive} & Generalizing DRL routing & Network
simulation & Class C topology/generalization evidence; no direct kernel ratio\\
Bi et al. \cite{bi2023dynamic} & Dynamic key-resource routing & Software and
network model & Class C resource-aware motivation\\
MadQCI \cite{martin2024madqci} & SDN-QKD deployment & Operational testbed & Class D
deployment context\\
Venkatachalam et al. \cite{venkatachalam2023postprocessing} & QKD post-processing &
FPGA implementation & Class D proof of FPGA suitability for a different function\\
This work & Fixed-point candidate evaluation & RTL, post-route, and physical replay &
Class A internal same-board comparison with exact reference agreement\\
\bottomrule
\end{tabularx}
\end{table}

The research position is consequently specific. \QFlow\ contributes physically
validated fixed-point candidate-evaluation evidence, while several software studies
contribute broader topology, algorithmic, or generalization evidence. Neither type
of evidence dominates the other in every dimension.

\section{Data-Analysis Principles}

The evaluation uses exact integer comparison for reference/RTL/board agreement;
post-route reports for maximum frequency and resources; accepted-request-to-done
cycles for kernel latency; profile counts, transitions, and dwell lengths for
adaptive behavior; and paired confidence intervals, sign tests, and effect sizes for
held-out route-quality comparisons. Missing power, energy, or CPU results are
treated as unmeasured rather than as zero.

\chapterconclusion{%
The background motivates key-resource-aware routing, establishes the distinction
between hard eligibility and soft profile-conditioned preference, and explains how
offline tabular Q-learning becomes deterministic policy-ROM inference. The
literature taxonomy positions \QFlow\ as a physically validated candidate evaluator
without constructing invalid speed comparisons to network algorithms or different
QKD components.}
